\scanpage{14}
(par exemple courbe ellipttque) est tel que $ {}_p A_K$ soit un groupe constant, alors est-il vrai que la trace $\handwritten{\phi}_E$ soit \deleted{nulle}\inserted{égale à $n=\dim A_K$} sur les \deleted{points}-éléments premiers à $p$ de $I$ ? (Seul le c\deleted{\unsure{q}}as de réduction non semi-stable de $A$ offre un problème; donc il faudrait ici trouver un exemple où le critère de réduction semi-stable de Raynaud est en défaut pour $\handwritten{\ell}=p$ !).

Dans le cas d'une courbe elliptique, et en car. rés. $\handwritten{\ne}2,\handwritten{3}$, prenant pour $K\handwritten{\unsure{'|(}}$ la plus petite extension sur laquelle on a réduction semi-stable, on trouve que $G$ est un sous-groupe de $\mathrm{Aut}(A_{s'})$, lequel \inserted{est cyclique} a\handwritten{(}$2,4$ ou $6$ élément\unsure{s}, donc est premier à $p$, et la question revient à celle de savoir si $\handwritten{\phi}_E$ est de valeur\deleted{\unsure{s}} constante égale à $2$, \deleted{si}\inserted{\handwritten{lorsque}} $ {}_p A$ est con\deleted{\unsure{s}}stant\deleted{\unsure{e}} \handwritten{(}et \deleted{\unsure{à}}si $K$ contient les racines $p$.èmes de l'unité\inserted{\handwritten{(?)}}; il suffit de regarder le cas où la courbe elliptique réduite, après extension $K\handwritten{\longrightarrow}K'$, est d'invariant de Hasse nul\handwritten{)}\deleted{\unsure{;}}.

% slanted
\margin{\handwritten{\uline{faux par} \uline{Tate}}}
% archive sheet number: 13

